\section{مبانی نظری}

\subsection{دیستروفی عضلانی دوشن}

دیستروفی عضلانی دوشن
\footnote{\lr{Duchenne Muscular Dystrophy (DMD)}} 
یکی از شایع‌ترین بیماری‌های ژنتیکی عصبی-عضلانی 
\footnote{\lr{Neuromuscular}}
در کودکان است که به‌صورت وابسته به کروموزوم 
\lr{X}
به ارث می‌رسد و عمدتاً پسران را درگیر می‌کند \cite{Emery2002DMD}.
این بیماری در اثر جهش در ژن دیستروفین
\footnote{\lr{Dystrophin}}
ایجاد می‌شود و در نتیجه آن، پروتئین دیستروفین که نقش مهمی در حفظ پایداری ساختاری سلول‌های عضلانی دارد، به مقدار کافی تولید نمی‌شود \cite{Hoffman1987Dystrophin}. 
نبود این پروتئین موجب آسیب تدریجی فیبرهای عضلانی و در نهایت جایگزینی آن‌ها با بافت چربی و همبند می‌شود. علاوه بر این دیستروفین نقش مهمی در پایداری مکانیکی غشای سلول‌های عضلانی دارد و اختلال در آن موجب آسیب‌پذیری فیبرهای عضلانی در برابر تنش‌های مکانیکی ناشی از انقباض می‌گردد.

از منظر بالینی، نشانه‌های بیماری غالباً در سال‌های ابتدایی کودکی ظاهر می‌شوند و می‌توانند شامل تأخیر در مهارت‌های حرکتی، دشواری در دویدن و بالا رفتن از پله‌ها، الگوی راه رفتن غیرطبیعی و زمین خوردن‌های مکرر باشند \cite{Bushby2010Diagnosis}. یکی از نشانه‌های شناخته‌شده در بسیاری از بیماران، 
\lr{Gowers' sign}
\footnote{الگوی جبرانی برخاستن از زمین} 
است که به ضعف عضلات کمربند لگنی نسبت داده می‌شود.

با پیشرفت بیماری، افت عملکرد حرکتی به‌صورت تدریجی رخ می‌دهد و در بسیاری از بیماران در سنین نوجوانی از دست رفتن توانایی راه رفتن مستقل مشاهده می‌شود. همچنین در مراحل پیشرفته‌تر، درگیری سیستم تنفسی و قلبی می‌تواند بر کیفیت زندگی و پیامدهای بالینی اثرگذار باشد \cite{Eagle2002Survival}.

\subsection{مقیاس \lr{NSAA}}

مقیاس
 \lr{NSAA}
 \footnote{\lr{North Star Ambulatory Assessment}}
 یکی از ابزارهای استاندارد ارزیابی عملکرد حرکتی در بیماران سرپایی مبتلا به دیستروفی عضلانی دوشن است که به‌طور گسترده در کارآزمایی‌های بالینی و مطالعات پایش طولی مورد استفاده قرار می‌گیرد \cite{Scott2009NSAA}. این مقیاس به‌طور اختصاصی برای بیماران دیستروفی عضلانی دوشن طراحی شده و هدف آن اندازه‌گیری توانایی انجام فعالیت‌های عملکردی روزمره مرتبط با راه رفتن و تعادل است.

\lr{NSAA}
 شامل ۱۷ فعالیت عملکردی است که طیفی از حرکات پایه تا فعالیت‌های دینامیک‌تر مانند پریدن و دویدن را دربرمی‌گیرد. هر فعالیت بر اساس یک مقیاس ترتیبی سه‌سطحی امتیازدهی می‌شود:

\begin{itemize}
	\item نمره 2: اجرای طبیعی و مستقل حرکت بدون استفاده از الگوهای جبرانی؛
	\item نمره 1: اجرای مستقل همراه با جبران حرکتی یا تغییر در الگوی طبیعی حرکت؛
	\item نمره 0: ناتوانی در انجام مستقل فعالیت.
\end{itemize}

حداکثر امتیاز قابل کسب در این مقیاس ۳۴ است که نشان‌دهنده عملکرد کامل حرکتی است. از آنجا که این مقیاس ماهیتی ترتیبی دارد، فاصله بین نمرات لزوماً بیانگر اختلاف عددی خطی در کیفیت حرکت نیست، بلکه نشان‌دهنده طبقه‌بندی سطوح عملکردی است.

\subsection{حسگرهای اینرسیایی}

حسگرهای اینرسیایی 
\footnote{ \lr{Inertial Measurement Unit (IMU)} }
 ابزارهایی پوشیدنی برای اندازه‌گیری پارامترهای سینماتیکی حرکت هستند. این واحدها عموماً شامل شتاب‌سنج
 \footnote{\lr{Accelerometer}}
  و ژیروسکوپ
 \footnote{\lr{Gyroscope}}
   سه‌محوره بوده و امکان ثبت شتاب خطی و سرعت زاویه‌ای را در راستای سه محور مختصاتی فراهم می‌کنند \cite{Sabatini2011IMU}.

شتاب‌سنج میزان شتاب خطی وارد بر حسگر را اندازه‌گیری می‌کند که شامل مؤلفه ناشی از حرکت و مؤلفه گرانشی است. بنابراین، در تحلیل حرکت لازم است اثر شتاب گرانشی از داده‌ها تفکیک یا جبران شود تا مؤلفه‌های دینامیکی حرکت استخراج گردند \cite{Madgwick2011SensorFusion}. ژیروسکوپ نیز سرعت زاویه‌ای حول محورهای مختصاتی را ثبت می‌کند و امکان تحلیل چرخش اندام‌ها و تغییرات وضعیت زاویه‌ای را فراهم می‌سازد.

داده‌های حاصل از حسگر اینرسیایی ممکن است تحت تأثیر نویز اندازه‌گیری، خطای کالیبراسیون و رانش انتگرالی\footnote{افزایش تدریجی خطا در تخمین زاویه که ناشی از انتگرال‌گیری طولانی‌مدت از سرعت زاویه‌ای است.} قرار گیرند. این موضوع به‌ویژه در کاربردهایی که نیاز به تخمین وضعیت زاویه‌ای دقیق دارند، اهمیت بیشتری پیدا می‌کند. ازاین‌رو، استفاده از روش‌های مناسب پیش‌پردازش و فیلترگذاری برای افزایش دقت و پایداری داده‌ها ضروری است.

\subsection{پیش پردازش سیگنال}

داده‌های خام حاصل از حسگرهای اینرسیایی معمولاً شامل نویز اندازه‌گیری، مؤلفه‌های ناخواسته و تغییرات غیرمرتبط با فعالیت هدف هستند. بنابراین پیش از انجام هرگونه تحلیل کمی، لازم است فرآیند پیش‌پردازش سیگنال انجام شود. فیلترگذاری مناسب برای حذف نویز و افزایش نسبت سیگنال به نویز یکی از گام‌های مهم در تحلیل داده‌ها است\cite{Winter2009Biomechanics}.

یکی از روش‌های رایج در تحلیل حرکت، استفاده از فیلترهای دیجیتال پایین‌گذر مانند فیلتر
 \lr{Butterworth}
 است. این فیلتر به دلیل پاسخ فرکانسی یکنواخت در باند عبور و عدم ایجاد نوسانات شدید در سیگنال، در کاربردهای بیومکانیکی متداول است
  \cite{Winter2009Biomechanics}.
  انتخاب فرکانس قطع مناسب باید با توجه به ماهیت حرکت و نرخ نمونه‌برداری انجام شود.

قطعه‌بندی زمانی
\footnote{ \lr{Segmentation}}
 نیز نقش مهمی در تحلیل فعالیت‌های حرکتی دارد. در این فرآیند، سیگنال پیوسته به بازه‌هایی متناظر با هر فعالیت مشخص تقسیم می‌شود تا امکان استخراج ویژگی‌های مرتبط با هر حرکت فراهم گردد. همچنین نرمال‌سازی
 \footnote{\lr{Data Normalization}}
  زمانی به مقایسه بین آزمودنی‌ها یا اجرای‌های مختلف یک فعالیت کمک می‌کند.

\subsection{استخراج ویژگی}

پس از انجام مراحل پیش‌پردازش، داده‌های حاصل از حسگرهای اینرسیایی همچنان به‌صورت سیگنال‌های خام و دنباله‌های زمانی در اختیار هستند. اگرچه این سیگنال‌ها اطلاعات کاملی از نحوه انجام حرکت را در خود دارند، اما استفاده مستقیم از آن‌ها در بسیاری از الگوریتم‌های یادگیری ماشین به دلیل ابعاد بالا، وجود هم‌بستگی میان نمونه‌ها و حساسیت به نویز، معمولاً منجر به افزایش هزینه محاسباتی و کاهش قابلیت تعمیم مدل می‌شود. ازاین‌رو، در بسیاری از مطالعات تحلیل حرکت، پیش از آموزش مدل، مجموعه‌ای از ویژگی‌های توصیف‌کننده از سیگنال استخراج می‌شود \cite{Bulling2014HumanActivity}.

استخراج ویژگی
\footnote{\lr{Feature Extraction}}
فرآیندی است که طی آن، اطلاعات مهم موجود در سیگنال خام به مجموعه‌ای از شاخص‌های عددی با ابعاد کمتر تبدیل می‌شود. هدف از این مرحله، حفظ اطلاعات مؤثر بر رفتار سیگنال و حذف اطلاعات غیرضروری یا تکراری است. ویژگی‌های استخراج‌شده می‌توانند نمایانگر خصوصیات آماری، زمانی یا فرکانسی سیگنال بوده و به الگوریتم یادگیری ماشین در شناسایی الگوهای موجود در داده کمک کنند.

به طور کلی، ویژگی‌های مورد استفاده در تحلیل سیگنال‌های حرکتی را می‌توان به دو گروه اصلی تقسیم کرد: ویژگی‌های حوزه زمان
\footnote{\lr{Time-domain Features}}
و ویژگی‌های حوزه فرکانس
\footnote{\lr{Frequency-domain Features}}.
ویژگی‌های حوزه زمان مستقیماً از مقادیر نمونه‌های سیگنال محاسبه می‌شوند و شاخص‌هایی مانند میانگین، انحراف معیار، واریانس، دامنه تغییرات، انرژی و چولگی را شامل می‌شوند. این ویژگی‌ها علاوه بر سادگی محاسبه، اطلاعات مناسبی از شدت حرکت، میزان تغییرات و رفتار کلی سیگنال ارائه می‌دهند.

در مقابل، ویژگی‌های حوزه فرکانس با تبدیل سیگنال به حوزه فرکانسی، معمولاً با استفاده از تبدیل فوریه
\footnote{\lr{Fast Fourier Transform (FFT)}}
استخراج می‌شوند. این ویژگی‌ها اطلاعاتی درباره توزیع انرژی سیگنال در فرکانس‌های مختلف، فرکانس غالب و سایر مشخصات طیفی ارائه می‌کنند و در برخی کاربردهای تحلیل حرکت و شناسایی فعالیت‌های انسانی مورد استفاده قرار می‌گیرند \cite{Bulling2014HumanActivity}.

\subsection{کاهش بعد داده‌ها}

در بسیاری از مسائل تحلیل داده، تعداد ویژگی‌های استخراج‌شده از سیگنال‌ها بسیار زیاد است. افزایش تعداد ویژگی‌ها علاوه بر افزایش زمان آموزش مدل‌ها، می‌تواند موجب افزایش پیچیدگی مدل و کاهش قابلیت تعمیم آن شود. به همین دلیل، از روش‌های کاهش بعد برای نمایش داده‌ها در فضایی با ابعاد کمتر استفاده می‌شود.

یکی از متداول‌ترین روش‌های کاهش بعد، تحلیل مؤلفه‌های اصلی
\footnote{\lr{Principal Component Analysis (PCA)}}
است. این روش با ایجاد ترکیب‌های خطی جدید از ویژگی‌های اولیه، داده‌ها را به مجموعه‌ای از مؤلفه‌های مستقل تبدیل می‌کند که بیشترین میزان واریانس موجود در داده را حفظ می‌کنند
 \cite{Jolliffe2002PCA}.


\subsection{یادگیری ماشین}

یادگیری ماشین 
\footnote{\lr{Machine Learning}}
 مجموعه‌ای از روش‌های محاسباتی است که به سیستم‌ها امکان می‌دهد الگوهای موجود در داده‌ها را یاد گرفته و بر اساس آن پیش‌بینی یا طبقه‌بندی انجام دهند 
 \cite{Bishop2006PatternRecognition}. 
 در تحلیل حرکت انسان، این روش‌ها به‌طور گسترده برای شناسایی فعالیت‌ها، طبقه‌بندی الگوهای حرکتی و پیش‌بینی شاخص‌های عملکردی به کار گرفته شده‌اند.

در چارچوب یادگیری نظارت‌شده
 \footnote{\lr{Supervised Learning}}
  مدل با استفاده از داده‌های برچسب‌دار
 \footnote{\lr{Labeled Data}}
   آموزش می‌بیند؛ به این معنا که برای هر نمونه داده، خروجی صحیح (مانند نوع فعالیت یا نمره عملکردی) مشخص است. پس از آموزش، مدل می‌تواند برای داده‌های جدید، خروجی متناظر را پیش‌بینی کند. روش‌هایی مانند ماشین بردار پشتیبان
 \footnote{\lr{Support Vector Machine}}
    درخت تصمیم
 \footnote{\lr{Decision Tree}}  
     و شبکه‌های عصبی مصنوعی 
 \footnote{ \lr{Artificial Neural Networks (ANN)}}
     از جمله الگوریتم‌هایی هستند که در مطالعات تحلیل فعالیت انسانی کاربرد گسترده‌ای داشته‌اند 
 \cite{Lara2013Survey}.

در زمینه تحلیل داده‌های حسگرهای پوشیدنی، یادگیری ماشین به‌عنوان ابزاری برای استخراج الگوهای پیچیده و غیرخطی از داده‌های چندبعدی شناخته می‌شود. ویژگی‌های استخراج‌شده از سیگنال‌های شتاب و سرعت زاویه‌ای می‌توانند به‌عنوان ورودی مدل قرار گرفته و رابطه آن‌ها با شاخص‌های عملکردی مدل‌سازی شود.

\subsection{معیارهای ارزیابی مدل}

برای ارزیابی عملکرد الگوریتم‌های یادگیری ماشین از معیارهای متداول طبقه‌بندی استفاده می‌شود. دقت
\footnote{\lr{Accuracy}}
نسبت نمونه‌های طبقه‌بندی‌شده صحیح به کل نمونه‌ها را نشان می‌دهد. با این حال، در مسائل دارای عدم توازن داده، استفاده از معیارهایی مانند
\lr{Precision}،
\lr{Recall}
و
\lr{F1-score}
اطلاعات دقیق‌تری از عملکرد مدل ارائه می‌دهد.

همچنین ماتریس درهم‌ریختگی
\footnote{\lr{Confusion Matrix}}
ابزاری مناسب برای تحلیل نحوه طبقه‌بندی نمونه‌ها و شناسایی خطاهای مدل محسوب می‌شود. در این پژوهش از تمامی این معیارها برای مقایسه عملکرد الگوریتم‌های مختلف استفاده شده است.

\subsection{ماشین بردار پشتیبان}

ماشین بردار پشتیبان
\footnote{\lr{Support Vector Machine (SVM)}}
یکی از پرکاربردترین الگوریتم‌های یادگیری نظارت‌شده برای مسائل طبقه‌بندی است. ایده اصلی این روش، یافتن ابرصفحه‌ای است که بتواند کلاس‌های مختلف را با بیشترین حاشیه ممکن از یکدیگر جدا کند \cite{Cortes1995SVM}.

در مواردی که داده‌ها به‌صورت خطی قابل جداسازی نباشند، با استفاده از توابع هسته
\footnote{\lr{Kernel Function}}
می‌توان داده‌ها را به فضای با ابعاد بالاتر نگاشت و جداسازی مناسب‌تری انجام داد. از متداول‌ترین هسته‌ها می‌توان به هسته خطی، چندجمله‌ای و تابع پایه شعاعی
\footnote{\lr{Radial Basis Function (RBF)}}
اشاره کرد.

\lr{SVM}
 توانایی بالا در مدل‌سازی مرزهای تصمیم پیچیده و عملکرد مناسب در مجموعه‌داده‌های با ابعاد زیاد دارد.

\subsection{الگوریتم نزدیک‌ترین همسایه}

الگوریتم نزدیک‌ترین همسایه
\footnote{\lr{K-Nearest Neighbors (KNN)}}
یکی از ساده‌ترین روش‌های طبقه‌بندی مبتنی بر نمونه است. در این روش، نمونه جدید بر اساس نزدیک‌ترین نمونه‌های موجود در مجموعه آموزش طبقه‌بندی می‌شود \cite{Cover1967KNN}.

ملاک نزدیکی معمولاً فاصله اقلیدسی است و پارامتر اصلی این الگوریتم، تعداد همسایه‌های مورد استفاده
($k$)
است. انتخاب مقدار مناسب برای این پارامتر تأثیر قابل توجهی بر عملکرد مدل دارد؛ به‌گونه‌ای که مقادیر کوچک ممکن است موجب حساسیت زیاد به نویز شوند و مقادیر بزرگ احتمال هموار شدن بیش از حد مرز تصمیم را افزایش دهند.


\subsection{جنگل تصادفی}

جنگل تصادفی
\footnote{\lr{Random Forest}}
یکی از روش‌های یادگیری مبتنی بر اجتماع مدل‌ها
\footnote{\lr{Ensemble Learning}}
است که از مجموعه‌ای از درخت‌های تصمیم برای انجام طبقه‌بندی استفاده می‌کند \cite{Breiman2001RandomForest}.

در این روش، هر درخت تصمیم با استفاده از زیرمجموعه‌ای تصادفی از داده‌های آموزشی و همچنین زیرمجموعه‌ای از ویژگی‌ها ساخته می‌شود. در مرحله پیش‌بینی، خروجی نهایی از طریق رأی‌گیری میان تمامی درخت‌ها تعیین می‌شود.

استفاده از چندین درخت تصمیم موجب کاهش بیش‌برازش، افزایش پایداری مدل و بهبود دقت پیش‌بینی می‌شود. به همین دلیل، 
\lr{Random Forest} 
در بسیاری از مسائل تحلیل داده‌های زیستی و داده‌های حاصل از حسگرهای پوشیدنی عملکرد مناسبی از خود نشان داده است.